\documentclass{includes/Thesis}

\usepackage{enumitem}
\usepackage{amssymb}

% Список первого уровня (itemize)
\setlist[itemize,1]{
    label=\textbullet,  % Символ: маленькая точка (точка по центру строки)
    leftmargin=1.75cm,          % Общий отступ слева для всего элемента списка
    itemsep=0\baselineskip,   % Вертикальный отступ между пунктами (можно 0pt для плотного списка)
    topsep=0\baselineskip     % Вертикальный отступ до и после всего списка
}

% Список второго уровня (itemize)
\setlist[itemize,2]{
    label=$\circ$,
    leftmargin=1.75cm,           % Общий отступ слева для всего элемента списка
    itemsep=0\baselineskip,
    topsep=0\baselineskip
}

% Список третьего уровня (itemize)
\setlist[itemize,3]{
    label=\scalebox{0.6}{$\blacksquare$},
    leftmargin=1.75cm,           % Общий отступ слева для всего элемента списка
    itemsep=0\baselineskip,
    topsep=0\baselineskip
}

\sloppy

\addbibresource{parts/refs.bib}

\AtBeginDocument{
    \renewcommand{\contentsname}{СОДЕРЖАНИЕ}
}

\begin{thesis}[
        institute = ФЕДЕРАЛЬНОЕ ГОСУДАРСТВЕННОЕ БЮДЖЕТНОЕ ОБРАЗОВАТЕЛЬНОЕ УЧРЕЖДЕНИЕ ВЫСШЕГО ОБРАЗОВАНИЯ «МОСКОВСКИЙ АВИАЦИОННЫЙ ИНСТИТУТ (НАЦИОНАЛЬНЫЙ ИССЛЕДОВАТЕЛЬСКИЙ УНИВЕРСИТЕТ)»,
        faculty = Компьютерные науки и прикладная математика,
        %
        title = Проектирование и реализация персонального ассистента для автоматизации управления задачами с использованием искусственного интеллекта,
        year = 2025,
        %
        authorGroup = М8О-411Б-22,
        educationalProgram = 02.03.02 «Фундаментальная информатика и информационные технологии»,
        authorName = М. С. Губарев,
        %
        academicTeacherTitle = {Кандидат физико-математических наук, доцент},
        academicTeacherName = А. М. Романенков,
        %
        academicSupervisorTitle = {Директор инстстута Компьютерные науки и прикладная математика, кандидат физико-математических наук, доцент},
        academicSupervisorName = С. С. Крылов,
        %
        isAcademic,
        UDC = 004.89,
        %
        keywordsRu={Curve25519; Ed25519; .NET; C\#; SIMD; AVX2; Криптография; Zero Allocation}
        keywordsEn = AI Assistant; Automated Task Management; Getting Things Done (GTD); Microservice Architecture; C++; YandexGPT; Telegram Bot,
    ]

    \setAbstractResource{parts/abstract-ru}{parts/abstract-en}

    \setTerminologyResource{parts/terminology}

    \setIntroResource{parts/intro}

    \addChapter{Теоретические основы и математический аппарат криптографических примитивов}{parts/chapter1}
    \addChapter{Проектирование архитектуры и программных решений высокопроизводительной криптографической библиотеки}{parts/chapter2}
    \addChapter{Практическая реализация и анализ производительности криптографических примитивов}{parts/chapter3}
    \addChapter{АНАЛИЗ РЕЗУЛЬТАТОВ И ПЕРСПЕКТИВЫ РАЗВИТИЯ}{parts/chapter4}

    \setConclusionResource{parts/conclusion}

    \addAppendix{Листинг кода}{parts/appendix-example-1}

\end{thesis}
